\documentclass[10pt,a4paper]{amsart}
\usepackage{amsmath,amssymb,amsthm,graphicx}
\usepackage{verbatim}
\newenvironment{reference}{\comment}{\endcomment}
%\newenvironment{reference}{}{}
\newenvironment{slogan}{\comment}{\endcomment}
\newenvironment{history}{\comment}{\endcomment}
% For commutative diagrams we use Xy-pic
\usepackage[all]{xy}

% We use 2cell for 2-commutative diagrams.
\xyoption{2cell}
% Theorem environments.
%
\theoremstyle{plain}
\newtheorem{theorem}[subsection]{Theorem}
\newtheorem{proposition}[subsection]{Proposition}
\newtheorem{lemma}[subsection]{Lemma}

\theoremstyle{definition}
\newtheorem{definition}[subsection]{Definition}
\newtheorem{example}[subsection]{Example}
\newtheorem{exercise}[subsection]{Exercise}
\newtheorem{situation}[subsection]{Situation}

\theoremstyle{remark}
\newtheorem{remark}[subsection]{Remark}
\newtheorem{remarks}[subsection]{Remarks}

\numberwithin{equation}{subsection}

% Macros
%
\def\lim{\mathop{\mathrm{lim}}\nolimits}
\def\colim{\mathop{\mathrm{colim}}\nolimits}
\def\Spec{\mathop{\mathrm{Spec}}}
\def\Hom{\mathop{\mathrm{Hom}}\nolimits}
\def\Ext{\mathop{\mathrm{Ext}}\nolimits}
\def\SheafHom{\mathop{\mathcal{H}\!\mathit{om}}\nolimits}
\def\SheafExt{\mathop{\mathcal{E}\!\mathit{xt}}\nolimits}
\def\Sch{\mathit{Sch}}
\def\Mor{\mathop{\mathrm{Mor}}\nolimits}
\def\Ob{\mathop{\mathrm{Ob}}\nolimits}
\def\Sh{\mathop{\mathit{Sh}}\nolimits}
\def\NL{\mathop{N\!L}\nolimits}
\def\CH{\mathop{\mathrm{CH}}\nolimits}
\def\proetale{{pro\text{-}\acute{e}tale}}
\def\etale{{\acute{e}tale}}
\def\QCoh{\mathit{QCoh}}
\def\Ker{\mathop{\mathrm{Ker}}}
\def\Im{\mathop{\mathrm{Im}}}
\def\Coker{\mathop{\mathrm{Coker}}}
\def\Coim{\mathop{\mathrm{Coim}}}

% Boxtimes
%
\DeclareMathSymbol{\boxtimes}{\mathbin}{AMSa}{"02}

%
% Macros for moduli stacks/spaces
%
\def\QCohstack{\mathcal{QC}\!\mathit{oh}}
\def\Cohstack{\mathcal{C}\!\mathit{oh}}
\def\Spacesstack{\mathcal{S}\!\mathit{paces}}
\def\Quotfunctor{\mathrm{Quot}}
\def\Hilbfunctor{\mathrm{Hilb}}
\def\Curvesstack{\mathcal{C}\!\mathit{urves}}
\def\Polarizedstack{\mathcal{P}\!\mathit{olarized}}
\def\Complexesstack{\mathcal{C}\!\mathit{omplexes}}
% \Pic is the operator that assigns to X its picard group, usage \Pic(X)
% \Picardstack_{X/B} denotes the Picard stack of X over B
% \Picardfunctor_{X/B} denotes the Picard functor of X over B
\def\Pic{\mathop{\mathrm{Pic}}\nolimits}
\def\Picardstack{\mathcal{P}\!\mathit{ic}}
\def\Picardfunctor{\mathrm{Pic}}
\def\Deformationcategory{\mathcal{D}\!\mathit{ef}}
\usepackage{cleveref}
\crefname{originalreference}{original reference}{original references}
\Crefname{originalreference}{Original reference}{Original references}
\makeatletter
\newlabel{categories-section-2-fibre-products}{{categories-section-2-fibre-products}{source}{}{sourcepointer-1}{}}
\newlabel{categories-section-2-fibre-products@cref}{{[originalreference][0][]categories-section-2-fibre-products}{[1][0][]source}{}{}{}}
\makeatother
\crefname{originallemma}{original lemma}{original lemmas}
\Crefname{originallemma}{Original lemma}{Original lemmas}
\makeatletter
\newlabel{lemma-base-change-affine}{{lemma-base-change-affine}{source}{}{sourcepointer-2}{}}
\newlabel{lemma-base-change-affine@cref}{{[originallemma][0][]lemma-base-change-affine}{[1][0][]source}{}{}{}}
\makeatother
\crefname{originallemma}{original lemma}{original lemmas}
\Crefname{originallemma}{Original lemma}{Original lemmas}
\makeatletter
\newlabel{lemma-composition-affine}{{lemma-composition-affine}{source}{}{sourcepointer-3}{}}
\newlabel{lemma-composition-affine@cref}{{[originallemma][0][]lemma-composition-affine}{[1][0][]source}{}{}{}}
\makeatother
\title{In a commutative triangle, affineness of one leg and of the other diagonal gives affineness of the third morphism}
\author{The Stacks Project Authors}
\date{stacks-project:stacks-morphisms.tex:local-snapshot-20261012; GNU FDL 1.2 or later}
\begin{document}
\maketitle
\noindent\textit{Editorial note.} All delivered mathematical text is the author's own, copied from the pinned Stacks Project chapter source sources/stacks-algebraic-stacks-morphisms/pinned-original-source.tex; no statement or proof step was rewritten, reordered or removed.\par\medskip
\noindent\textit{Editorial note.} Preamble provenance: the Stacks Project keeps one single shared preamble for every chapter; this card takes it directly from the staged sources/stacks-algebraic-stacks-morphisms/preamble.tex as three verbatim spans (lines 8-12, 14-18 and 154-217), so no derived file was written and nothing was added or redefined. The wrapper class and the base packages are supplied by the assembler.\par\medskip
\noindent\textit{Editorial note.} Bibliography provenance: neither the counted statement nor its proof contains a citation (this lane rescans the whole span before writing the card), so the card carries no bibliography and no entry was reconstructed.\par\medskip
\noindent\textit{Editorial note.} Selection provenance: this card was opened under the file-based control contract (CONTROL.md), with the mechanical pre-card gate run on its statement and proof interval before the card was written.\par\medskip
\noindent\textit{Allowed standard prerequisites (editorial).}\begin{itemize}
\item The base change permanence of affineness. Invoked twice inside the counted proof (source lines 2374-2376 and 2379-2382). Makes the projection of the fibre product affine and the graph morphism affine.
\item The identification of the graph with a base change of the diagonal of the second leg. Used inside the counted proof (source lines 2377-2381). Realizes the graph morphism as a base change of that diagonal.
\item The composition permanence of affineness. Invoked at the end of the counted proof (source lines 2383-2385). Combines the two affine factors into the assertion.
\end{itemize}\medskip
\ifdefined\hypertarget\hypertarget{sourcepointer-1}{}\fi
\noindent\textit{Source reference.} Original reference categories-section-2-fibre-products: sources/stacks-algebraic-stacks-morphisms/pinned-original-source.tex, source lines 2382--2382. Navigation pointer only: the counted text refers to the Stacks item categories-section-2-fibre-products, which belongs to another chapter of the project and is not part of this card. The pointer carries the target's own label as its number because this chapter's staged snapshot does not contain the other chapter's tag table; its content is not reproduced here.\par\medskip
\ifdefined\hypertarget\hypertarget{sourcepointer-2}{}\fi
\noindent\textit{Source reference.} Original lemma lemma-base-change-affine: sources/stacks-algebraic-stacks-morphisms/pinned-original-source.tex, source lines 2376--2376. Navigation pointer only: the counted text refers to the Stacks item lemma-base-change-affine, which belongs to another chapter of the project and is not part of this card. The pointer carries the target's own label as its number because this chapter's staged snapshot does not contain the other chapter's tag table; its content is not reproduced here.\par\medskip
\ifdefined\hypertarget\hypertarget{sourcepointer-3}{}\fi
\noindent\textit{Source reference.} Original lemma lemma-composition-affine: sources/stacks-algebraic-stacks-morphisms/pinned-original-source.tex, source lines 2386--2386. Navigation pointer only: the counted text refers to the Stacks item lemma-composition-affine, which belongs to another chapter of the project and is not part of this card. The pointer carries the target's own label as its number because this chapter's staged snapshot does not contain the other chapter's tag table; its content is not reproduced here.\par\medskip
\begin{lemma}
\label{lemma-affine-permanence}
Let
$$
\xymatrix{
\mathcal{X} \ar[rr]_f \ar[rd]_a & & \mathcal{Y} \ar[dl]^b \\
& \mathcal{Z}
}
$$
be a commutative diagram of morphisms of algebraic stacks.
If $a$ is affine and $\Delta_b$ is affine, then
$f$ is affine.
\end{lemma}

\begin{proof}
The base change
$\text{pr}_2 : \mathcal{X} \times_\mathcal{Z} \mathcal{Y} \to \mathcal{Y}$
of $a$ is affine by Lemma \ref{lemma-base-change-affine}. The morphism
$(1, f) : \mathcal{X} \to \mathcal{X} \times_\mathcal{Z} \mathcal{Y}$
is the base change of
$\Delta_b : \mathcal{Y} \to \mathcal{Y} \times_\mathcal{Z} \mathcal{Y}$
by the morphism $\mathcal{X} \times_\mathcal{Z} \mathcal{Y} \to
\mathcal{Y} \times_\mathcal{Z} \mathcal{Y}$ (see material in
Categories, Section \ref{categories-section-2-fibre-products}).
Hence it is affine by
Lemma \ref{lemma-base-change-affine}.
The composition $f = \text{pr}_2 \circ (1, f)$ of affine morphisms is affine by
Lemma \ref{lemma-composition-affine} and the proof is done.
\end{proof}

\end{document}
